\chapter{The Development Plan}
\label{ch:development-plan}

Chapter~\ref{ch:complete-pipeline} gives the logical order for layering features and models.
This chapter gives the actual build order: what to implement first given project priorities and foundation work that must happen before anything else.
Priority ordering: trading signal $>$ academic rigor $>$ model novelty.


%% ──────────────────────────────────────────────
\section{Milestones}
\label{sec:milestones}
%% ──────────────────────────────────────────────

Table~\ref{tab:milestones} defines seven milestones with acceptance criteria and dependencies.

\begin{table}[htbp]
\centering
\caption{Development milestones. M1--M5 form the minimum viable deliverable.}
\label{tab:milestones}
\small
\begin{tabular}{@{}clp{4.8cm}p{4.2cm}l@{}}
\toprule
\textbf{M\#} & \textbf{Milestone} & \textbf{Acceptance Criteria} & \textbf{Key Tasks} & \textbf{Deps} \\
\midrule
1 & Fix Foundation &
  Purge gap $\geq h$ enforced; $\QLIKE$ sign matches \citet{Patton2011}; context kwarg added; all existing tests pass &
  CV fix, $\QLIKE$ fix, protocol extension, safe\_log dedup &
  --- \\[8pt]
2 & LightGBM &
  Custom $\QLIKE$ objective converges; Optuna finds improved params; walk-forward OOS predictions for 3 horizons &
  $\QLIKE$ gradient/hessian, model class, Optuna, walk-forward &
  M1 \\[8pt]
3 & Tournament &
  8 models $\times$ 3 horizons table with DM $p$-values on dev universe &
  Run baselines, DM test, tournament table &
  M2 \\[8pt]
4 & Layer~2 Options &
  Options features produce daily values with no look-ahead; $\QLIKE$ lift documented &
  OptionsLayer, IV surface wiring, validation &
  M1 \\[8pt]
5 & Signal &
  IV--RV gap signal; equity curve; positive OOS Sharpe &
  Signal logic, P\&L backtest, performance metrics &
  M3, M4 \\[8pt]
6 & Ensemble &
  Residual stacking and prediction blending tested; 10-model tournament table &
  Residual stacking, inverse-$\QLIKE$ blending &
  M3 \\[8pt]
7 & Stretch &
  Ordered by impact: regime $\QLIKE$, MCS, LSTM feature, full universe, figures, Rashomon &
  Each task independent &
  M3--M6 \\
\bottomrule
\end{tabular}
\end{table}


%% ──────────────────────────────────────────────
\section{Critical Path}
\label{sec:critical-path}
%% ──────────────────────────────────────────────

Figure~\ref{fig:critical-path} shows milestone dependencies.
The critical path (M1 $\to$ M2 $\to$ M3 $\to$ M6 $\to$ M7) is highlighted.
M4 branches from M1 and runs in parallel with M2--M3.
M5 and M6 can run in parallel after M3 completes.

\begin{figure}[htbp]
\centering
\resizebox{\textwidth}{!}{%
\begin{tikzpicture}[
  node distance=0.8cm and 1.6cm,
  >=Stealth,
  milestone/.style={draw, rounded corners, minimum height=0.9cm, minimum width=2.4cm,
    align=center, font=\small, fill=blue!8},
  critical/.style={milestone, line width=1.5pt, fill=orange!12},
  arrow/.style={-{Stealth[length=2.5mm]}, thick},
  critarrow/.style={arrow, line width=1.5pt, color=defblue},
]
  % Critical path
  \node[critical] (m1) {M1\\Foundation};
  \node[critical, right=of m1] (m2) {M2\\LightGBM};
  \node[critical, right=of m2] (m3) {M3\\Tournament};
  \node[critical, right=of m3] (m6) {M6\\Ensemble};
  \node[critical, right=of m6] (m7) {M7\\Stretch};

  \draw[critarrow] (m1) -- (m2);
  \draw[critarrow] (m2) -- (m3);
  \draw[critarrow] (m3) -- (m6);
  \draw[critarrow] (m6) -- (m7);

  % Parallel track
  \node[milestone, below=1.2cm of m2] (m4) {M4\\Options};
  \node[milestone, below=1.2cm of m6] (m5) {M5\\Signal};

  \draw[arrow] (m1) -- (m4);
  \draw[arrow] (m4) -| (m5);
  \draw[arrow] (m3) |- (m5);
  \draw[arrow] (m5) -| (m7);
\end{tikzpicture}%
}
\caption{Milestone dependencies.
Bold path: critical path (M1--M2--M3--M6--M7).
M4 runs in parallel with M2--M3.
M5 requires both M3 (forecasts) and M4 (options features).
M7 waits on all prior milestones.}
\label{fig:critical-path}
\end{figure}


%% ──────────────────────────────────────────────
%% Boxes
%% ──────────────────────────────────────────────

\begin{keyidea}[M1--M5 Is the Minimum Viable Deliverable]
A $\QLIKE$ tournament (7~HAR variants + LightGBM, DM tests) plus a tradeable IV--RV signal with P\&L backtest is a presentable result.
Everything in M6--M7 is upside.
If time runs out after M5, the project has a defensible outcome.
\end{keyidea}

\begin{warning}[M1 Is Non-Negotiable]
The purge gap bug causes silent data leakage for $h=22$.
The $\QLIKE$ sign convention determines whether the loss function penalizes over-prediction or under-prediction correctly.
All results produced before M1 is complete are potentially invalid.
Do not skip ahead.
\end{warning}
